\PassOptionsToPackage{xetex}{xcolor}
\documentclass[shortabstract,english,mgr]{iithesis}

\usepackage[utf8]{inputenc}
\usepackage{unicode-math}
\setmathfont{XITS Math}
\usepackage{src_tex/agda}

\polishtitle    {Weryfikacja kompilatora dla języka programowania z efektami sterowania}
\englishtitle   {Verified compiler for language with control effects}
\polishabstract {\ldots}
\englishabstract{\ldots}
\author         {Paweł Dybiec}
\advisor        {dr Piotr Polesiuk}
%\date          {}                     % Data zlozenia pracy
% Dane do oswiadczenia o autorskim wykonaniu
\transcriptnum {271900}                     % Numer indeksu
\polishabstract{Tematem tej pracy jest stworzenie i zweryfikowanie języka z ograniczonymi
                operatorami kontroli. Używa on polimorficznych typów aby pozwolić na wyrażenie
                obliczeń generycznych wobec pewnych efektów. Celem pracy jest zweryfikowanie poprawnosci
                definicji takiego języka oraz jego translacji do maszyny abstrakcyjnej.}
\englishabstract{We present type system for language with labeled delimited control operators.
                 It utilizes polymorphic types and effects to allow expressions parametrized by effects.
                 Our goal is to formally prove correctness of such language and present translation
                 of it to simple abstract machine.}
%\advisorgen    { } % Nazwisko promotora w dopelniaczu
%%%%%

%%%%% WLASNE DODATKOWE PAKIETY
%
%\usepackage{graphicx,listings,amsmath,amssymb,amsthm,amsfonts,tikz}
%
%
%\theoremstyle{definition} \newtheorem{definition}{Definition}[chapter]
%\theoremstyle{remark} \newtheorem{remark}[definition]{Observation}
%\theoremstyle{plain} \newtheorem{theorem}[definition]{Theorem}
%\theoremstyle{plain} \newtheorem{lemma}[definition]{Lemma}
%\renewcommand \qedsymbol {\ensuremath{\square}}
% ...
%%%%%
\renewcommand{\AgdaCodeStyle}{\tiny}
\usepackage{fontspec}
\setmonofont[Scale=0.9]{JuliaMono}
\newfontfamily{\AgdaMonoFont}{JuliaMono}
\newfontfamily{\AgdaTypewriterFont}{JuliaMono}
\renewcommand{\AgdaSpace}{\texttt{\AgdaMonoFont{} }}
\renewcommand{\AgdaFontStyle}[1]{{\AgdaTypewriterFont{}#1}}
\renewcommand{\AgdaKeywordFontStyle}[1]{{\AgdaMonoFont{}#1}}
\renewcommand{\AgdaStringFontStyle}[1]{{\AgdaMonoFont{}#1}}
\renewcommand{\AgdaCommentFontStyle}[1]{{\AgdaTypewriterFont{}#1}}
\renewcommand{\AgdaBoundFontStyle}[1]{\AgdaTypewriterFont{}#1}

\usepackage{hyperref}

\begin{document}


\chapter{Introduction}

Control operators have been present in languages for a long time, one noticable example is SCHEME programming language which started in 70s\cite{scheme}. Operators for delimited control introduced by Danvy, Filinski\cite{abstracting} and Felleisen\cite{prompts} allow to control them more precisely. They still allow for non-local execution, which enables intresting computational effects such as emulating non-determinism, exceptions, failure and others.

In this paper we try to build typed version of such calculus with polymorphic types (similiar to Piróg, Polesiuk and Sieczkowski\cite{typed_equivalence}) and named labels. Usually such calculi jump from shift to closest enclosing reset. Here we tried labeling them to allow stacking of different effects. This was planned originally to use approach similiar to described by Materzok, Biernacki\cite{subtyping} to compile language to abstract machine.

Formalisation work was done in Agda 2.8. Code archive has also nix derivation describing exact environment to build both package and this paper to ensure reproducibility.


\chapter{Types}
\input{src_tex/Types}
\chapter{Runtime}
\input{src_tex/Runtime}
\chapter{Progress}
\input{src_tex/Progress}
\chapter{Machine}
\input{src_tex/Machine}
\chapter{Discussion/Closing thoughts}
In this paper we presented language with labeled delimited effect handlers together with partial formalisation of such language. We defined 3 term languages, typing judgements for 2 of them, translations between those languages and that typing is preserved.
We defined intrinsically typed frames and metaframes that allow us to define reduction relation, how it preserves types. Since they are typed we statically know what effects are handled by frame.
We also defined draft of abstract machine and how it could evaluate this language.
This formalisation could be expanded and finished to fully prove correctness of the language and translations.

\begin{thebibliography}{1}
\bibitem{typed_equivalence} Maciej Piróg,
  Piotr Polesiuk,
  Filip Sieczkowski. \href{https://doi.org/10.4230/LIPIcs.FSCD.2019.30}{Typed Equivalence of Effect Handlers and
Delimited Control}, In 4th International Conference on Formal Structures for Computation and Deduction (FSCD 2019), pages 30:1-30:16, 2019
\bibitem{subtyping} Marek Materzok,
  Dariusz Biernacki. \href{https://dl.acm.org/doi/abs/10.1145/2034574.2034786}{Subtyping Delimited Continuations} In ACM SIGPLAN Notices, Volume 46, Issue 9, pages 81-93, 2011
\bibitem{abstracting} O. Danvy and A. Filinski. \href{https://dl.acm.org/doi/10.1145/91556.91622}{Abstracting control} In LISP and Functional Programming, pages 151-160, 1990
\bibitem{prompts} M. Felleisen. \href{https://dl.acm.org/doi/10.1145/73560.73576}{The theory and practice of first-class prompts} In POPL '88: Proceedings of the 15th ACM SIGPLAN-SIGACT symposium on Principles of programming languages, pages 180-190, 1988
\bibitem{scheme}  Guy Lewis Steele, Gerald Jay Sussman. \href{https://dspace.mit.edu/handle/1721.1/6283}{The Revised Report on SCHEME: A Dialect of LISP} in MIT Artificial Intelligence
Memo 452, 1978

\end{thebibliography}


\end{document}
